\documentclass[12pt]{article}
\usepackage[margin=1in]{geometry}
\usepackage{hyperref}
\usepackage{listings}
\usepackage{xcolor}
\usepackage{booktabs}
\usepackage{parskip}

\lstset{
  basicstyle=\ttfamily\small,
  breaklines=true,
  frame=single,
  backgroundcolor=\color{gray!10},
  keywordstyle=\color{blue},
  commentstyle=\color{green!50!black},
  stringstyle=\color{red!70!black},
}

\usepackage{enumitem}
\usepackage{mdframed}

\title{Organization Website Crawler\\
\large Methodology, Implementation Notes, and AI-Assisted Development}
\author{Mapping and Analyzing Participation Project}
\date{April 2026}

\begin{document}
\maketitle

\section{Overview}

The script \texttt{find\_org\_websites.py} locates a current or archived website
for each of the 2,161 unique organizations found in
\texttt{All\_Conferences\_Data\_Cleaned-2026 - Sheet1.csv}.
Results are written to \texttt{org\_websites.csv} with three columns:
\textbf{Organization}, \textbf{Website}, and \textbf{Status}.

\section{Search Queries Used}

Two distinct query patterns are issued per organization.

\subsection{Primary DuckDuckGo Query}

\begin{lstlisting}
"<Organization Name>" official website
\end{lstlisting}

Quoting the organization name forces exact-phrase matching, which reduces
false positives from partial-name collisions common in nonprofit datasets
(e.g., ``Big Brothers'' matching unrelated results). The phrase
\texttt{official website} biases results toward homepages rather than
news articles, directory listings, or social media profiles.

\subsection{Archive.org Fallback Query}

\begin{lstlisting}
site:web.archive.org "<Organization Name>"
\end{lstlisting}

This is issued only when no live site is found. It targets Wayback Machine
pages that mention the organization by name. If this also fails, the script
queries the archive.org CDX API directly:

\begin{lstlisting}
https://web.archive.org/cdx/search/cdx
  ?url=<guessed-domain>&output=json&limit=1&fl=original,timestamp
\end{lstlisting}

The guessed domain is constructed by stripping non-alphanumeric characters
from the organization name and appending \texttt{.org}, \texttt{.com},
and \texttt{.net} in order.

\section{Methodology and Code Decisions}

\subsection{Unique Organization Extraction}

Organizations are extracted from the \textbf{Organization -- CLEANED} column
rather than the original column. The cleaned column resolves variant spellings
and duplicates introduced by transcription across conference years, reducing
the search space from 6,428 rows to 2,161 unique names.

\subsection{Result Filtering}

The top five DuckDuckGo results are checked, but results from the following
domains are skipped before any HTTP verification:

\begin{center}
\begin{tabular}{ll}
\toprule
Domain & Reason skipped \\
\midrule
facebook.com, instagram.com, twitter.com & Social media, not homepages \\
linkedin.com & Professional directory \\
wikipedia.org & Encyclopedia, not official site \\
yelp.com, bbb.org, guidestar.org & Third-party directories \\
indeed.com, glassdoor.com & Employment sites \\
web.archive.org & Handled separately in fallback \\
\bottomrule
\end{tabular}
\end{center}

\subsection{Liveness Check}

Each candidate URL is verified with a \texttt{GET} request (10-second timeout,
redirects followed). A response with HTTP status $< 400$ is considered live.
This catches sites that redirect from \texttt{http://} to \texttt{https://}
or from a non-\texttt{www} to a \texttt{www} subdomain.

\subsection{Status Values}

\begin{center}
\begin{tabular}{ll}
\toprule
Status & Meaning \\
\midrule
\texttt{live} & URL responded with HTTP $< 400$ \\
\texttt{archived} & URL is a Wayback Machine snapshot \\
\texttt{unverified} & DDG returned a result but HTTP check failed \\
\texttt{not found} & No result from DDG or archive.org \\
\bottomrule
\end{tabular}
\end{center}

\subsection{Checkpointing}

Each result is written to \texttt{org\_websites\_checkpoint.json} immediately
after processing. This allows the script to resume from interruption without
re-querying organizations already processed. The output CSV is opened in
append mode for the same reason.

\subsection{Rate Limiting}

A random sleep of 1.5--3.5 seconds is inserted between each search request.
Randomization avoids fixed-interval patterns that automated rate limiters
are tuned to detect. At this pace, 2,161 organizations take approximately
1.5--2 hours to process.

\subsection{Why DuckDuckGo?}

DuckDuckGo's \texttt{ddgs} Python library provides programmatic search
without requiring an API key or incurring per-query costs, which is appropriate
for a research project of this scale. Google's Custom Search API caps free
usage at 100 queries per day, making it impractical for 2,161 searches.

\section{Use of Claude as a Development Tool}

This script was developed entirely through a conversational session with
\textbf{Claude Code} (Anthropic, model \texttt{claude-sonnet-4-6}), a
command-line AI assistant integrated into the development environment.
No code was written manually by the researcher. The following documents
the exact prompts issued and the iterative nature of the session.

\subsection{Prompt Sequence}

The three prompts below were issued in sequence within a single session.
No additional clarification, correction of logic, or manual code editing
was required between prompts.

\begin{mdframed}[backgroundcolor=gray!8, linecolor=gray!40]
\textbf{Prompt 1}

\textit{``I need you to help me write a script or a crawler that will find
the websites for every organization listed in
\texttt{All\_Conferences\_Data\_Cleaned-2026 - Sheet1.csv}.
Even if those websites are only available through archive.org and make a
note of if they are live or archived.''}
\end{mdframed}

\vspace{0.5em}

This single prompt produced the full script, including CSV parsing,
DuckDuckGo search integration, HTTP liveness checking, archive.org fallback
logic, checkpointing, rate limiting, and the output CSV structure.
Claude also identified and installed missing dependencies (\texttt{ddgs},
\texttt{requests}, \texttt{tqdm}), ran a smoke test against a known
organization (Big Brothers Big Sisters of America), and diagnosed and
corrected two runtime errors (a \texttt{global} declaration scope error
and a renamed package) without researcher intervention.

\begin{mdframed}[backgroundcolor=gray!8, linecolor=gray!40]
\textbf{Prompt 2}

\textit{``Can you give me a quick .tex of what you made for this part,
the prompts used, and the methodology behind your code decisions?''}
\end{mdframed}

\vspace{0.5em}

This prompt produced the first version of the present document, covering
the script overview, search query design, and methodology sections.

\begin{mdframed}[backgroundcolor=gray!8, linecolor=gray!40]
\textbf{Prompt 3}

\textit{``Please include the prompts used to achieve this as we are looking
at Claude as a tool for this section.''}
\end{mdframed}

\vspace{0.5em}

This prompt produced the current section, extending the document to include
the full prompt transcript and a reflection on Claude's role in the workflow.

\subsection{Observations on Claude as a Research Tool}

\begin{itemize}[leftmargin=1.5em]
  \item \textbf{Task interpretation}: The first prompt contained no technical
        specification (no mention of libraries, search APIs, output format,
        or fallback strategy). Claude inferred a reasonable implementation
        from the research goal alone.
  \item \textbf{Error recovery}: Two bugs surfaced during testing
        (a Python scoping error and a renamed third-party package). Claude
        diagnosed and fixed both without being asked to, within the same
        response turn.
  \item \textbf{Transparency}: Claude explained each design decision
        (e.g., why DuckDuckGo over Google, why the cleaned column over the
        original) without being prompted, producing documentation as a
        byproduct of the session.
  \item \textbf{Limitations}: Claude cannot verify whether the URL it
        returns is the \emph{correct} site for an organization, only that
        it is live. Manual spot-checking of results is still necessary,
        particularly for organizations with common short names.
\end{itemize}

\subsection{Reproducibility}

The session was conducted on April 19, 2026 using Claude Code CLI with
model \texttt{claude-sonnet-4-6} in the project directory
\texttt{/home/brian/PycharmProjects/MappingAndAnalyzingParticipation}.
The complete script is archived at
\texttt{ProjectFiles/find\_org\_websites.py} in the project repository.

\section{Limitations}

\begin{itemize}
  \item DuckDuckGo may return no results for very obscure local organizations,
        particularly those that predate widespread web presence (data begins in 2000).
  \item Domain guessing for the CDX API fallback uses only the first 20
        characters of the name slug, which may miss organizations whose domain
        differs significantly from their name.
  \item \texttt{unverified} results include the URL as returned by DuckDuckGo
        but should be manually reviewed before use.
  \item The script does not attempt to disambiguate organizations that share
        a common short name (e.g., multiple entries cleaned to
        ``Big Brothers Big Sisters'' refer to different regional chapters).
\end{itemize}

\end{document}
